%!TEX root = ../main.tex
\chapter*{Sommario}
\addcontentsline{toc}{chapter}{Sommario}

Un incrocio semaforico deve garantire un flusso regolare dei veicoli che
vi transitano. Le soluzioni classiche affrontano il problema attraverso
un piano a tempi fissi, con fasi di durata prestabilita calcolata su
stime di traffico medie, oppure attraverso un controllo attuato dai sensori, che
concede il verde a una direzione solo quando un sensore rileva la
presenza di un veicolo in coda.

Quando gli incroci da controllare sono più di uno, emerge un problema
ulteriore: le decisioni prese a un incrocio devono essere coordinate con
quelle degli incroci vicini, perché il traffico smaltito da un'intersezione si riversa in
parte in quello in entrata ad un'altra intersezione. Sistemi come SCOOT e SCATS rispondono a
questa esigenza coordinando i piani semaforici di un'intera area urbana
sulla base del traffico osservato in tempo reale.

Entrambe le famiglie di soluzioni si basano però su euristiche definite
manualmente e restano per questo poco flessibili di fronte a condizioni di
traffico diverse da quelle previste in fase di progettazione. Il
reinforcement learning offre un'alternativa, apprendendo una
politica di controllo dall'esperienza simulata invece di codificarla
esplicitamente. MetaSTGAT (Wang et al., 2022) ne è un esempio, con un agente RL per
intersezione coordinato da una \emph{graph attention network} i cui pesi
sono generati dinamicamente da un modulo di meta-learning condizionato su
caratteristiche locali della rete stradale. Il lavoro originale si
concentra però sull'efficienza aggregata della rete, senza verificare se
il controllore generalizzi a topologie mai osservate in addestramento né
se il vantaggio ottenuto sia distribuito equamente tra le direzioni di
traffico servite.

Questo lavoro replica da zero l'architettura e la procedura di training
di MetaSTGAT, estendendola esplicitamente su questi due assi, equità e
generalizzazione, su un ambiente costruito su \textsc{CityFlow} con reti
a griglia di tre dimensioni ($4\times4$, $5\times5$, $6\times6$) e
traffico generato parametricamente.
Ogni differenza introdotta rispetto alla procedura originale è isolata da un
sistema di preset di ablazione e motivata singolarmente.

I risultati sperimentali confrontano le baseline classiche (Fixed-Time,
MaxPressure) con tre varianti del modello proposto che pesano
diversamente nella ricompensa la componente di equità, con uno studio di
ablazione completo e con un confronto tra tre meccanismi di comunicazione
spaziale alternativi. Il termine di ricompensa dedicato all'equità
risulta essere utile per il raggiungimento della stessa e
il meccanismo di comunicazione spaziale scelto si conferma efficace
per l'ottenimento di un buon coordinamento tra gli incroci della rete.
